% ============================================================================
% aintec_additions.tex — paste-ready blocks for aintec_panel (AINTEC '26 draft)
% Each block states WHERE it goes. Figures needed in figures/:
%   ratio_cdf_all3.png, cdn_heatmap.png   (from experiments/07_longitudinal_panel/analysis/figures/)
% ============================================================================


% ----------------------------------------------------------------------------
% BLOCK 1 — add at the END of §3.1 Site Selection (after "…produces identical output.")
% ----------------------------------------------------------------------------
\paragraph{Location verification.}
Because the hosting classification above rests on geo-IP, we verify it against physics once
measurements are available: no reply can arrive faster than light in vacuum would need for the
round trip to a site's claimed location, so any measured RTT below that bound proves the claimed
location wrong. Applied to the first days of the panel, this test invalidated the registered
locations of 34 of the 40 CDN-class sites (anycast addresses whose registration city is, e.g.,
Toronto, yet which answer Pakistani probes in under 3\,ms) and corrected six unicast sites: three
sites classed as Pakistani-hosted were in fact far from every probe (nearest ping over 200\,ms,
consistent with their US and Singapore geo-IP) and are reclassified as Abroad, and three sites
whose claimed city was physically impossible were relocated by latency multilateration, using each
probe's RTT as a distance bound. The corrected hosting split is therefore 37 Pakistani, 40 CDN,
and 23 Abroad, and all per-class results below use the corrected classes. We treat this not as a
flaw in the sample but as a finding in its own right---among the misclassified sites is a
provincial government service hosted in the United States---and as independent validation that
registration-based location data cannot be trusted at exactly the addresses this study cares
about.


% ----------------------------------------------------------------------------
% BLOCK 2 — §3.2 Vantage Points: replace the sentence
%   "…at the launch of the week-long run reported here this comprised 17 probes."
% with the following text (and keep the rest of the subsection as is)
% ----------------------------------------------------------------------------
% ...at the launch of the week-long run reported here this comprised 17 probes, of which 16
% returned data (one probe, on the state operator NTC, was offline throughout). We additionally
% verified each probe's metadata against the platform API at analysis time rather than trusting
% run-time labels: one probe's stored label carried a wrong autonomous-system number and was
% corrected (it is a Z-Com vantage), and one probe declares a placeholder coordinate and is
% excluded from all distance-based analysis while its round-trip times, which remain valid, are
% retained elsewhere.


% ----------------------------------------------------------------------------
% BLOCK 3 — NEW subsection. Insert AFTER §3.3 Measurement Design,
% BEFORE §3.4 Tromboning Classification.
% ----------------------------------------------------------------------------
\subsection{The Latency Ratio}
\label{sec:ratio}
To compare connections fairly across distances, we normalise each measured round-trip time by the
fastest round trip physics allows for that connection, giving a dimensionless \emph{latency
ratio}: a value of $1$ means the path performs at its physical limit, and a value of $k$ means the
connection takes $k$ times longer than the best possible. Because the ratio is unitless it can be
pooled across near and far sites and plotted as a distribution, and we retain the hosting class of
every site so the distributions can be sliced by type. The theoretical minimum differs by class.

For \textbf{unicast sites} (the Pakistani and Abroad classes, one physical server each), the
minimum is set by geography. With the probe and server coordinates, the great-circle distance $d$
follows from the Haversine formula, and light in optical fibre covers it at
$v_f = c/1.468 \approx 204{,}200$\,km/s, so
\begin{equation}
  \mathit{RTT}_{\min} \;=\; \frac{2d}{v_f} \;\approx\; \frac{d~[\mathrm{km}]}{102}~\mathrm{ms},
  \qquad
  \text{ratio} \;=\; \frac{\mathit{RTT}_{\mathrm{measured}}}{\mathit{RTT}_{\min}},
\end{equation}
where $\mathit{RTT}_{\mathrm{measured}}$ is the minimum ping over the window (noise only ever adds
delay, so the minimum best approximates the true path latency). Probe coordinates come from the
platform API; server coordinates from the verified geo-IP of the measured address
(\S3.1). Pairs closer than 30\,km are excluded from the ratio---the denominator approaches zero
and the quotient degenerates---and are reported directly in milliseconds instead (their median in
our data is 6\,ms, i.e.\ same-city delivery is healthy).

For \textbf{CDN sites} no single server location exists: the same anycast address is served from
different points of presence for different ISPs, and its registered location is provably wrong
(\S3.1). We therefore define the theoretical minimum \emph{empirically}: the best RTT that any
Pakistani probe achieves to that site. This floor is attainable by construction---any ISP could
peer with the best-performing network, at the exchange this paper is about, and match its
path---so the ratio
\begin{equation}
  \text{ratio}_{\mathrm{CDN}} \;=\;
  \frac{\mathit{RTT}_{\mathrm{measured}}(\mathrm{probe}, s)}
       {\min_{\mathrm{all\ probes}} \mathit{RTT}_{\mathrm{measured}}(\cdot\,, s)}
\end{equation}
measures exactly what an ISP loses by not peering. It contains no distance and no geo-IP; the
best-connected ISP scores $1$ and defines the frontier.


% ----------------------------------------------------------------------------
% BLOCK 4 — §4 Results: REPLACE the current section intro paragraph
% ("The week-long run reported here is in collection…") with the following,
% then insert the two new subsections 4.1 and 4.2 BEFORE the existing RQ1
% subsection (which becomes 4.3; RQ2 -> 4.4; RQ3 -> 4.5, no text changes).
% ----------------------------------------------------------------------------
The week-long run is in collection at the time of writing; its window closes on 18~July~2026,
after which the traceroute and ping halves are merged on probe, site, and round, and the three
analyses the panel is designed for---one per research question---are computed on the full week
(\S\ref{sec:rq1}--\S\ref{sec:rq3} below). Ahead of that, we report two supporting analyses from a
snapshot taken three days into the window. They serve two purposes: they exercise the full
pipeline of \S3 end to end, and---because every number below is a minimum over rounds, which
additional days can only tighten---they establish stable baselines that the full-week results will
refine rather than overturn. Both are labelled preliminary.

\subsection{Preliminary: Latency Ratio by Site Type}
\label{sec:prelim-ratio}
Figure~\ref{fig:ratio-cdf} shows the distribution of the latency ratio (\S\ref{sec:ratio}) for the
507 probe--site connections to unicast sites with a ping response and the 617 connections to CDN
sites, as cumulative distribution functions. Three observations stand out. First, connections to
\emph{foreign} sites are the control group the metric needs: their ratios are tight (median
2.5$\times$, none beyond ${\sim}7\times$), showing what ordinary routing achieves when it works.
Second, connections to \emph{Pakistani} sites match that typical performance (median 2.9$\times$)
but roughly a third of them tail out to 10--30$\times$---inefficiency that, given the control
group, cannot be blamed on ordinary Internet overhead and sits exactly where domestic detours
live; the full-week traceroutes attribute these paths in \S\ref{sec:rq1}. Third, the \emph{CDN}
curve is bimodal: about a third of ISP--site pairs sit at ${\approx}1\times$ (those ISPs ride the
locally peered path and define the floor), while roughly a quarter exceed 40$\times$. Access to
anycast content is effectively binary---an ISP either peers locally or hauls the traffic to a
distant point of presence---which previews the cross-section of the next subsection. One caveat
qualifies the Pakistani curve: 22 of the 40 Pakistani-class sites drop ICMP and are absent here;
they will be filled in from the TCP traceroute half of the panel, whose final hop the web servers
must answer.

\begin{figure}[t]
  \centering
  \includegraphics[width=\columnwidth]{ratio_cdf_all3.png}
  \caption{Latency ratio by site type (snapshot, first three days). For Pakistani and Abroad sites
  the theoretical minimum is straight fibre over the verified great-circle distance; for CDN sites
  it is the best RTT any Pakistani probe achieves (\S\ref{sec:ratio}). Foreign sites are uniform;
  Pakistani sites carry a long tail; CDN access is bimodal.}
  \label{fig:ratio-cdf}
\end{figure}

\subsection{Preliminary: The CDN Cross-Section---the Local Path Exists}
\label{sec:prelim-cdn}
RQ3 asks whether a local path exists for traffic that nonetheless leaves the country, and answers
it longitudinally, through verdicts that flip over the week. The CDN measurements already answer
the same question \emph{cross-sectionally}. Figure~\ref{fig:cdn-heatmap} shows, for each ISP and
each CDN-hosted site, how close the point of presence that serves that ISP is: local
($<$15\,ms), regional (15--50\,ms), or far ($>$50\,ms), thresholds we calibrated in earlier
cache-presence work. The rows order themselves into a peering ranking: the best-peered network
reaches 85\% of CDN content locally at a median of 3\,ms, two others reach 41\% and 20\%, and the
remaining six ISPs---including the largest operator---reach essentially none of it locally,
with medians up to 136\,ms. Expressed with the CDN ratio of \S\ref{sec:ratio}, the median cost of
non-peering ranges from $2\times$ (second-best network) to $52\times$ (the largest operator): the
same national content, in some cases military and provincial-government sites, is tens of times
slower depending only on which ISP a citizen uses. The columns separate causes: sites red across
\emph{every} row are far because of their own offshore fronting (no peering can fix them), whereas
columns green for well-peered ISPs and red for the rest quantify precisely the local path that
exists and is not taken. This is the paper's thesis in cross-section; the longitudinal version
follows on the full week in \S\ref{sec:rq3}.

\begin{figure}[t]
  \centering
  \includegraphics[width=\columnwidth]{cdn_heatmap.png}
  \caption{CDN access grid (snapshot, best probe per ISP). Green: served from a local PoP
  ($<$15\,ms); amber: regional (15--50\,ms); red: far ($>$50\,ms). Rows form a peering ranking;
  solid-red columns are sites fronted offshore for every ISP.}
  \label{fig:cdn-heatmap}
\end{figure}

% ----------------------------------------------------------------------------
% BLOCK 5 — labels to add on the existing RQ subsections (for the cross-refs above):
%   \subsection{Extent and Attribution of Tromboning (RQ1)}\label{sec:rq1}
%   \subsection{The Latency Cost and Its Stability (RQ2)}\label{sec:rq2}
%   \subsection{Existence of a Local Path (RQ3)}\label{sec:rq3}
% Also in §4.4 (RQ2), add one sentence at the end:
%   "The preliminary ratio distributions of \S\ref{sec:prelim-ratio} are the partial-window
%    baseline this analysis refines: the full week adds the diurnal decomposition and the
%    per-ISP split, and a path-only variant that subtracts each probe's measured last-mile
%    floor (on the snapshot, the Pakistani median falls from 2.9x to 1.9x, below the foreign
%    2.4x, indicating the typical domestic route is efficient and the cost lives in the
%    access overhead and the tail)."
% ----------------------------------------------------------------------------
